\section*{Abstract}

Cuban agriculture faces a multiplicative crisis characterized by shortages of raw materials, energy, water, and soil degradation, exacerbated by climate change and economic blockade. Irrigation infrastructure has collapsed, with obsolete technologies and absence of precision agriculture, while institutional management proves deficient due to irrigation scheduling without water balances and lack of trained personnel. The intersection with the energy crisis—scheduled blackouts, fuel shortages, and lack of spare parts—paralyzes pumping systems, including the most efficient ones. Furthermore, progressive soil salinization from overexploitation without adequate drainage generates salt stress that reduces productivity.

This work proposes the technological assimilation of the Internet of Things (IoT) as an approach to mitigate poor water resource management in the Cuban agricultural context. The current environment of real-time data transmission technologies will be described, the necessary artifacts to generate a metrics system for resource control will be analyzed, and technology assimilation will be executed under conditions of economic and energy constraints. The study object focuses on data transmission tools in IoT environments applied to precision agriculture, seeking to contribute to food sovereignty and national security established in governmental policies until 2030.


%\begin{description}
%	\item[Key words:]{the list of words separated by commas (,)}
%\end{description}
%\end Abstract